\section{Design}\label{sec:design}

\sys has four parts in a chain (\cref{fig:overview}): an eBPF probe that turns each system call of a watched container
into a token and counts per-thread 4-token sequences in the kernel; a per-workload count language model, learnt only
from normal behaviour, that scores every 0.5\,s window; an alert path that pushes an over-threshold window straight to
the cluster's Alertmanager; and a root-cause analysis (RCA) that builds a process evidence tree around each delivered
alert. The kernel only builds tokens and counts grams; all probabilities are computed in user space. A model is learnt
offline, locked by its SHA-256 and unchanged during a run.

\begin{figure}[t]
\centering
\fbox{\parbox{0.95\columnwidth}{\centering\small \pending{figure: system overview (translated from the thesis)}}}
\caption{\sys on one worker and on the master; the fit runs offline on normal traffic.}
\label{fig:overview}
\end{figure}

\subsection{Tokens that carry return values}\label{sec:tokens}

Many exploits keep the set of system calls an application makes and change only its data: a path outside the usual
directories, an unusual read or write size, an error code. A token therefore keeps, besides the system-call name, one
piece of the return information, chosen by priority:
\begin{enumerate}
\item \emph{Error code}: if the call is one of 27 that keep their errno and it failed, the token is
  \texttt{name!errno}, e.g.\ \texttt{openat!2} for \texttt{ENOENT}. The 27 are the stat family, \texttt{access} and
  \texttt{faccessat}, the open family, \texttt{read}, \texttt{write}, \texttt{writev}, \texttt{ioctl},
  \texttt{connect}, \texttt{execve}, \texttt{fcntl}, \texttt{close}, \texttt{rename}, \texttt{renameat},
  \texttt{sendfile}, \texttt{unlink}, \texttt{unlinkat}, \texttt{mkdir}, \texttt{kill}, \texttt{recvfrom},
  \texttt{shutdown} and \texttt{futex}.
\item \emph{Successful open}: \texttt{name:mode:dir} with mode \texttt{r} or \texttt{w} and the first component of the
  absolute path (\texttt{/etc}, \texttt{/tmp}; \texttt{/} at the root, \texttt{rel} for a relative path).
\item \emph{Successful read or write}: \texttt{name\#b} with $b=\min(15,\lfloor\log_2 n\rfloor)$ for $n$ bytes
  returned, e.g.\ \texttt{read\#12} for 4\,KiB.
\item Otherwise the name alone.
\end{enumerate}
In the kernel a token is a 32-bit number: system-call number in bits 23--31, token kind in bits 21--22, and in the rest
the error code, the size bucket, or a write flag with an FNV-1a hash of the first directory. The list of 27 calls and
the size resolution were chosen on the development part of LID-DS~2021 with attack labels; \cref{sec:eval-design}
shows that the set is wider than needed and which parts carry the signal.

\subsection{Per-thread 4-grams counted in the kernel}

Ordered steps such as open, read, close only mean something within one thread, and a service interleaves many threads.
Each thread therefore keeps its last four tokens, and event $e_i$ yields the 4-gram $g_i=(t_{i-3},t_{i-2},t_{i-1},t_i)$
of that thread (shorter at thread start). The thread history is not cut at window borders. The gram is counted in the
0.5\,s window that holds the event, in a kernel map keyed by (cgroup, window, gram); single events never cross to user
space.

\subsection{A count language model}\label{sec:model}

The model estimates $P(t_i\mid t_{i-3},t_{i-2},t_{i-1})$ by counting normal behaviour; it has no neural network and no
gradient training. Let $c(h,w)$ be the number of times token $w$ followed context $h$, $c(h)=\sum_w c(h,w)$, and $T(h)$
the number of distinct tokens that followed $h$.

\paragraph{Base level.} With $\nu(w)$ the system-call name of $w$, $N$ the number of events, $V$ the number of distinct
names and $V_\nu$ the number of distinct tokens of name $\nu$,
\begin{equation}
P_0(w)=\frac{c(\nu)+\frac12}{N+\frac12(V+1)}\cdot\frac{c(\nu,w)+\frac12}{c(\nu)+\frac12(V_\nu+1)} .
\end{equation}
The second factor is the probability of the variant (error code, size, directory) given the name. If the name itself
was never seen, it is replaced by the mean probability of an unseen variant over the known names, so an unseen
\texttt{openat!13} stays far more likely than a never-called \texttt{ptrace}.

\paragraph{Interpolated Witten--Bell smoothing}~\cite{witten1991,chen1999} goes from short to long contexts. With $h_k$
the $k$ tokens before $w$, $k=1,2,3$,
\begin{equation}
P_k(w\mid h_k)=\frac{c(h_k,w)+T(h_k)\,P_{k-1}(w\mid h_{k-1})}{c(h_k)+T(h_k)},
\end{equation}
and $P_k=P_{k-1}$ if $h_k$ was never seen; the score uses $P_3$. A context that is always followed by the same few
tokens trusts its own counts; one with many continuations hands more mass to the shorter context. A new 4-gram made of
familiar steps therefore gets a moderate probability, while one containing a strange token gets a tiny one.

\paragraph{Window score.} The surprise of an event is $-\ln P_3$, and the score of window $W$ is the mean of its ten
largest surprises, each gram counted with its multiplicity:
\begin{equation}
s(W)=\tfrac{1}{10}\textstyle\sum_{\text{top 10 } e\in W}-\ln P_3(e),
\end{equation}
the mean of all events if $W$ has fewer than ten, and $0$ if it is empty. The top of the window keeps a few strange
events from being diluted by thousands of familiar ones; ten rather than one keeps a single rare normal event from
raising an alarm. As an example from a live run, a harmless control binary in a RabbitMQ container produced a window of
89 events in which the gram (\texttt{getpid}, \texttt{clock\_nanosleep}, \texttt{getpid}, \texttt{clock\_nanosleep}),
absent from RabbitMQ's normal behaviour, occurred six times with surprise 34.0 and the next gram six times with 7.0:
$s=(6\cdot 34.0+4\cdot 7.0)/10=23.2$, above the workload's threshold of 13.4. The same binary in another service
produced the gram only three times, $s=11.3$, below that workload's 13.3. The model measures novelty and its density,
not intent.

\subsection{A threshold from normal behaviour only}\label{sec:threshold}

In each unit of normal behaviour (a recording of a dataset, or a 60\,s block of one container on the cluster) take
the largest window score $M_j$. Scoring normal data with the model that learnt it would understate the scores of new
data, so the scores are held out: the learning data is split into five folds and each fold is scored by the model of
the other four. On a dataset a recording goes to fold $\mathrm{sha256}(\mathit{name})\bmod 5$; on the cluster, time is
cut into contiguous 30-minute chunks dealt to the folds in turn, so a held-out window is far from the windows that were
learnt. The model that scores live data learns all of it. We drop the largest 1\,\% of the maxima (at least one), fit a
Gumbel distribution~\cite{coles2001} by moments,
$\beta=\frac{\sqrt6}{\pi}\sigma_M$, $\mu=\bar M-0.5772\,\beta$, and set
\begin{equation}
\tau=\mu-\beta\ln\bigl(-\ln(1-\alpha')\bigr),\qquad \alpha'=0.0002\times0.9 .
\end{equation}
Equivalently $\tau=\bar M+6.27\,\sigma_M$. The Gumbel assumption does not hold
(\cref{sec:eval-design}): $\alpha'$ is a knob, not a guaranteed false-alarm rate, and false alarms are measured.
A window with $s(W)>\tau$ raises an alarm at the end of that window; one window suffices because many exploits show in
a single window.

\subsection{Alert path and root-cause analysis}

A bridge on the master reads each worker's decisions from its local disk and pushes every alarm to Alertmanager at once,
instead of exposing it to a periodic Prometheus scrape (\cref{sec:eval-alert} measures both). For each alert that
Alertmanager received, the RCA opens a case and builds a process evidence tree from Tetragon events (process, fork and
flow policies) around the alarm, ranks the new executions near it as suspects (ties broken by the launch time closest
to the alarm), and merges later alarms of the same container within a cool-down into the same case.

\subsection{Lineage: the workload's own processes versus operator sessions}\label{sec:lineage}

The model flags almost every new process that enters a container through \texttt{kubectl exec}, harmless or not
(\cref{sec:eval-benign}). A workload model answers ``did this container do something it never did?'', and an operator
session always does. We therefore tag every thread in the kernel as \emph{native}, part of the workload's own process
tree, or \emph{session}, entered through a runtime exec, together with all its descendants. At fork
(\texttt{sched\_process\_fork}) the child copies its parent's tag, so a child cannot change lineage by outliving its
parent. A thread first seen without state takes the tag of its thread-group leader, else of its parent if the parent is
in the same cgroup, else it is native when its pid in its own pid namespace is~1 (the container's entrypoint) and
session otherwise, because a runtime exec joins the namespace from a parent outside the cgroup. The tag is part of the
gram key, so the loader reports each window unchanged together with its session part. In the latest run the tag is
telemetry only: the alert path still scores the whole window.
